\section{Positive-$\epsilon$ Stale-Subspace Failure}

We study the quadratic
\[
f(x,y)=\frac{1}{2}\left(a x^2 + b y^2\right),\qquad a>0,\; b>0,
\]
under a fixed stale subspace $P=e_1$. The projected branch sees only the
$x$-coordinate, while the orthogonal branch coincides with the $y$-gradient.
With the positive-$\epsilon$ scalar recovery used in our corrected theorem
regime,
\[
R_t = a x_t,\qquad
\psi_t(R_t)=\frac{R_t}{|R_t|+\epsilon_a},\qquad
\phi_t = \frac{|\psi_t(R_t)|}{|R_t|+\epsilon_s}
= \frac{|a x_t|}{(|a x_t|+\epsilon_a)(|a x_t|+\epsilon_s)},
\]
the updates are
\[
x_{t+1}=x_t-\eta \frac{a x_t}{|a x_t|+\epsilon_a},
\qquad
y_{t+1}=y_t-\eta \phi_t b y_t.
\]

\begin{theorem}[Positive-$\epsilon$ scalar recovery can leave a nonzero orthogonal gradient]
\label{thm:positive-epsilon-failure}
Consider the dynamics above with $0<\eta a<2\epsilon_a$ and
\[
M_0 = \max_{0\le r\le a|x_0|}\frac{r}{(r+\epsilon_a)(r+\epsilon_s)},
\qquad
\eta b M_0 \le \frac{1}{2}.
\]
Then there exists a nonempty set of initial conditions for which
\[
x_t \to 0,\qquad \sum_{t=0}^\infty \phi_t < \infty,
\]
while
\[
y_t \to y_\infty \neq 0,\qquad \norm{\grad f(x_t,y_t)} \to b|y_\infty| > 0.
\]
\end{theorem}

The interpretation is direct. Positive-$\epsilon$ normalization can stabilize
the projected coordinate and still send only finite cumulative mass through the
orthogonal recovery branch. The theorem does not require projection failure in
the sense of never observing the full gradient. Instead, the failure is caused
by the recovery scale itself.

The corrected theorem-regime diagnostic in Figure~\ref{fig:theorem-regime}
verifies exactly this pattern: $\phi_t$ decays, $\sum_t \phi_t$ plateaus, and
the raw recovery branch leaves a nonzero final $y_t$ and gradient norm.
